\documentclass{article}
\usepackage{graphicx}
\usepackage{booktabs}
\usepackage{hyperref}
\usepackage{amsmath}
\usepackage{geometry}
\usepackage{listings}
\usepackage{xcolor}
\geometry{margin=1in}

\lstset{
  basicstyle=\ttfamily\small,
  breaklines=true,
  frame=single,
  backgroundcolor=\color{gray!10},
}

\title{EE/CS 148B HW1: Build a Language Model from Scratch}
\author{Alex Chung}
\date{May 2026}

\begin{document}

\maketitle
\tableofcontents
\newpage

% ─────────────────────────────────────────────────────────────────────────────
\section{Tokenizer}
% ─────────────────────────────────────────────────────────────────────────────

\subsection{Insight (BPE): Naive Merging}

Naively merging pairs of bytes across raw text can create separate tokens for
words that are semantically almost identical but differ only in punctuation.
For example, \texttt{dog!} and \texttt{dog?} could become completely different
tokens, even though both contain the same word \texttt{dog} and should share
much of their meaning.  This matters because downstream models must then learn
that \texttt{dog!} and \texttt{dog?} are related concepts from context alone,
wasting model capacity that could have been used for genuinely distinct
vocabulary.

\subsection{BPE Training on TinyStories}

Training the byte-level BPE tokenizer on TinyStories completed in 6.7 minutes,
which is well within the 30-minute limit.  The final tokenizer had
10{,}000 tokens total, consisting of 256 initial byte tokens, 1 special token
(\texttt{<|endoftext|>}), and 9{,}743 learned merges.

The longest token in the vocabulary was \texttt{ accomplishment}, which is
15 bytes long.  This result makes sense because \texttt{accomplishment} is a
long English word that appears often enough in the TinyStories dataset for BPE
to merge its bytes into a single token.  The leading space is also expected
because the GPT-2-style pre-tokenizer treats space-prefixed words as units,
so the tokenizer correctly learns frequent subword patterns in their natural
word-boundary context.

% ─────────────────────────────────────────────────────────────────────────────
\section{Transformer Architecture}
% ─────────────────────────────────────────────────────────────────────────────

The model is a decoder-only Transformer language model built entirely from
primitive operations (no \texttt{torch.nn.functional} calls).  Key components:

\begin{itemize}
  \item \textbf{Embedding}: token embedding lookup + sinusoidal positional
        embedding (fixed, non-learnable).
  \item \textbf{Pre-norm blocks}: LayerNorm $\to$ Multi-head Causal Attention
        $\to$ residual; LayerNorm $\to$ FFN (SiLU + linear) $\to$ residual.
  \item \textbf{Attention}: scaled dot-product with lower-triangular causal
        mask; numerically stable softmax (subtract row-wise max before
        \texttt{exp}).
  \item \textbf{Output}: final LayerNorm, then a linear projection to
        vocabulary logits.
\end{itemize}

The baseline model configuration is shown in Table~\ref{tab:config}.

\begin{table}[h]
\centering
\caption{Baseline model hyperparameters}
\label{tab:config}
\begin{tabular}{ll}
\toprule
Hyperparameter & Value \\
\midrule
Vocabulary size & 10{,}000 \\
Context length & 256 tokens \\
Model dimension $d_{\text{model}}$ & 512 \\
Number of layers & 4 \\
Number of heads & 8 \\
Feed-forward dimension $d_{\text{ff}}$ & 2{,}048 \\
Total parameters & 22{,}832{,}128 \\
\bottomrule
\end{tabular}
\end{table}

% ─────────────────────────────────────────────────────────────────────────────
\section{Training}
% ─────────────────────────────────────────────────────────────────────────────

\subsection{Training Setup}

\begin{itemize}
  \item \textbf{Optimizer}: AdamW ($\beta_1=0.9$, $\beta_2=0.95$,
        $\varepsilon=10^{-8}$, weight decay $= 0.1$).
  \item \textbf{LR schedule}: linear warmup for 200 steps, then cosine
        annealing from $3\times10^{-4}$ to $3\times10^{-5}$.
  \item \textbf{Batch size}: 64 sequences of 256 tokens each
        ($\approx 16{,}384$ tokens per step).
  \item \textbf{Gradient clipping}: max norm 1.0.
  \item \textbf{Data}: TinyStories training split ($\approx541$M tokens,
        memory-mapped).
  \item \textbf{Device}: Apple M-series GPU (MPS backend).
\end{itemize}

\subsection{Reaching Validation Loss $\leq 2.0$}

Phase 1 ran for 5{,}000 steps at $\text{lr}=3\times10^{-4}$, achieving a
best validation loss of $2.1036$ (perplexity $= 8.20$).  The cosine schedule
had reached its minimum ($3\times10^{-5}$) by the end of phase 1, leaving no
room for further improvement with the same optimizer state.

To cross the $2.0$ threshold, a \emph{warm restart} was performed: the model
weights were loaded from the best phase-1 checkpoint, the AdamW optimizer was
reset to a fresh state, and training continued for 2{,}000 more steps at
$\text{lr}=5\times10^{-4}$ (warmup 100 steps, min-lr $5\times10^{-5}$).
This allowed the optimizer to escape the flat region around the local minimum
found by the cosine schedule.

The model crossed validation loss $2.0$ at step 1{,}000 of the warm restart
and reached a final best validation loss of $\mathbf{1.8452}$
(perplexity $\mathbf{6.33}$) at step 2{,}000.

\begin{figure}[h]
  \centering
  \includegraphics[width=0.85\textwidth]{plots/baseline_curve.png}
  \caption{Baseline training curve.  Phase 1 (lr=$3\times10^{-4}$, 5{,}000
    steps) reaches val\_loss $=2.10$.  The warm-restart phase (lr=$5\times10^{-4}$,
    fresh optimizer, 2{,}000 steps) crosses the target of $2.0$ at
    step $\approx1{,}000$ and finishes at $1.845$.}
  \label{fig:baseline}
\end{figure}

% ─────────────────────────────────────────────────────────────────────────────
\section{Text Generation}
% ─────────────────────────────────────────────────────────────────────────────

\subsection{Sampling Implementation}

Text is generated autoregressively.  At each step:
\begin{enumerate}
  \item The last $\min(\text{len}(\text{tokens}),256)$ token IDs are fed to
        the model to obtain logits for the next token.
  \item \textbf{Temperature scaling}: logits are divided by $\tau$, then
        numerically-stable softmax is applied:
        \[
          p_i = \frac{\exp((l_i - \max_j l_j)/\tau)}
                     {\sum_k \exp((l_k - \max_j l_j)/\tau)}
        \]
  \item \textbf{Top-$p$ (nucleus) filtering}: tokens are sorted descending by
        probability; the smallest prefix whose cumulative probability $\geq p$
        is retained, probabilities are renormalized, and a token is sampled
        with \texttt{torch.multinomial}.
  \item Generation halts on \texttt{<|endoftext|>} or after
        \texttt{max\_tokens} new tokens.
\end{enumerate}

\subsection{Insight: Temperature Scaling}

\textbf{When $\tau \to 0$}: the highest-logit token dominates
exponentially; in the limit this is equivalent to greedy argmax decoding ---
the model always picks its single most likely next token, producing repetitive
but locally coherent text.

\textbf{When $\tau \to \infty$}: all logits are divided by a huge number,
making them nearly equal before softmax, so the distribution approaches
uniform --- the model samples uniformly at random from the entire vocabulary,
producing incoherent noise.

Values in between (e.g.\ $\tau = 0.7$--$1.0$) trade off between diversity
and quality.

\subsection{Generated Text Sample}

The following 190-token sample was generated from the best checkpoint
(val\_loss $= 1.845$) using temperature $\tau = 0.8$ and top-$p = 0.92$:

\vspace{0.5em}
\begin{lstlisting}
Prompt: Once upon a time, there was a little girl named Lily who loved to
explore the forest near her home.

Generated:
One day, while playing in her home, she saw a big red ball near her house.
She was very curious about the ball.
Lily wanted to see the big red ball. She walked up to the big red ball and
started to follow it. As she walked, she saw a big, red ball. She wanted to
play with the big red ball, but she was scared.
Lily tried to push the big red ball, but it was too fast. She thought of a
plan to stop the ball. She took the big red ball and ran to the big red ball.
Lily saw the big red ball and felt happy.
But then, something unexpected happened. The big red ball started to glow!
It was a magic ball. Lily and the magic ball did not get the ball back. Lily
was so surprised and happy. They played with the big red ball all day, and
the big red ball never came back.
<|endoftext|>
\end{lstlisting}

\subsubsection*{Fluency Commentary}

The output is \textbf{grammatically correct} throughout and conforms to the
``TinyStories'' register: simple past tense, short sentences, a clear
protagonist and a problem-resolution arc.  However, the text has noticeable
weaknesses:

\begin{enumerate}
  \item \textbf{Repetition / broken coreference}: The phrase ``big red ball''
        appears nine times in 190 tokens.  The model has learned the surface
        co-occurrence between the noun phrase and the story genre but lacks a
        mechanism to track entities and vary referring expressions.

  \item \textbf{Logical inconsistency}: Lily ``ran to the big red ball'' after
        already holding it, and the magic ball ``never came back'' after she
        played with it all day.  Transformer language models predict the next
        token conditioned on a fixed context window and do not maintain a
        persistent world state, so they cannot guarantee global narrative
        consistency.
\end{enumerate}

Two factors that most strongly affect output quality:
\begin{itemize}
  \item \textbf{Model perplexity / training budget}: with more steps or a
        larger model the distribution over next tokens is sharper and more
        contextually appropriate, directly reducing repetition and incoherence.
  \item \textbf{Context length}: the 256-token context window limits how far
        back the model can track entities and plot; a longer context or a
        memory mechanism would improve narrative consistency.
\end{itemize}

% ─────────────────────────────────────────────────────────────────────────────
\section{Insight: Perplexity}
% ─────────────────────────────────────────────────────────────────────────────

Perplexity $\text{PPL} = e^{\mathcal{L}}$ (where $\mathcal{L}$ is the
average cross-entropy loss per token) is preferred over raw loss for two
reasons:

\begin{enumerate}
  \item \textbf{Interpretable units}: perplexity can be understood as the
        ``effective branching factor'' --- the model is, on average, as
        uncertain as if it had to choose uniformly among $\text{PPL}$ equally
        likely next tokens.  A perplexity of 6.33 means the model is, on
        average, confused between roughly 6 options per token, which is an
        intuitive measure of language understanding.

  \item \textbf{Tokenizer-independent comparisons}: cross-entropy loss
        depends on the tokenization granularity (character-level models have
        lower per-token loss than word-level models simply because characters
        are easier to predict).  Perplexity normalises for this by measuring
        uncertainty per token in the model's own vocabulary, making results
        more comparable across systems with different tokenizers.
\end{enumerate}

% ─────────────────────────────────────────────────────────────────────────────
\section{Ablation 1: Effect of LayerNorm}
% ─────────────────────────────────────────────────────────────────────────────

\subsection{Setup}

Two ``no-LayerNorm'' variants were trained: one at the same learning rate as
the baseline ($3\times10^{-4}$) to test for instability, and one at a lower
rate ($3\times10^{-5}$) to see whether stability can be recovered with a
reduced learning rate.  Both ran on the same architecture (4 layers, 512-dim)
with all LayerNorm modules removed.

\begin{table}[h]
\centering
\caption{LayerNorm ablation results}
\label{tab:ln}
\begin{tabular}{lrrrr}
\toprule
Run & Steps & Best val loss & Best PPL & LR \\
\midrule
Baseline (with LN) & 5{,}000 + 2{,}000 & 1.845 & 6.33 & $3\times10^{-4}$ \\
No-LN, lr=$3\times10^{-4}$ & 400 & 3.534 & 34.26 & $3\times10^{-4}$ \\
No-LN, lr=$3\times10^{-5}$ & 1{,}500 & 3.948 & 51.85 & $3\times10^{-5}$ \\
\bottomrule
\end{tabular}
\end{table}

\begin{figure}[h]
  \centering
  \includegraphics[width=0.85\textwidth]{plots/layernorm_ablation.png}
  \caption{LayerNorm ablation learning curves.  Both no-LayerNorm variants
    converge far more slowly than the baseline and reach much higher
    validation loss.  Lowering the learning rate stabilises training
    superficially but actually \emph{worsens} final performance.}
  \label{fig:ln}
\end{figure}

\subsection{Observations and Explanation}

\textbf{Why the high-LR no-LN run degrades:}
Without LayerNorm, the residual stream activations can grow unboundedly.  The
attention softmax is sensitive to the scale of its inputs (large dot products
push softmax toward saturation, collapsing the attention distribution to
one-hot), and the FFN nonlinearity similarly saturates for large pre-activations.
At lr $= 3\times10^{-4}$ these effects compound over layers: gradients arrive
with very different magnitudes across layers, causing erratic parameter updates
and slow convergence.  The model stalls at val\_loss $\approx 3.5$ after only
400 steps.

\textbf{Why the low-LR no-LN run is even worse:}
Intuitively, lowering the learning rate reduces the gradient signal but does
\emph{not} fix the underlying problem: activations are still unscaled and the
optimization landscape is still poorly conditioned.  The model therefore makes
smaller, noisier updates while still suffering from the same saturation issues,
leading to slower convergence and a worse final loss ($3.95$ after 1{,}500 steps).
LayerNorm's benefit is not merely ``keeping gradients small'' but ensuring that
each layer receives inputs with roughly unit variance, which re-centers the
optimization landscape so that gradient descent steps are actually useful.

\textbf{Conclusion:}  LayerNorm is critical for training stability and fast
convergence in Transformers.  Removing it degrades performance dramatically
regardless of the learning rate, because the root cause is ill-conditioned
layer activations, not learning-rate magnitude alone.

% ─────────────────────────────────────────────────────────────────────────────
\section{Ablation 2: Sinusoidal PE vs.\ NoPE}
% ─────────────────────────────────────────────────────────────────────────────

\subsection{Setup}

The NoPE variant removes positional embeddings entirely: the input to the
Transformer is the token embedding alone.  All other hyperparameters are
identical to the baseline (4 layers, 512-dim, lr $= 3\times10^{-4}$,
2{,}500 steps).

\begin{table}[h]
\centering
\caption{Positional embedding ablation results at 2{,}500 steps}
\label{tab:pe}
\begin{tabular}{lrr}
\toprule
Model & Val loss & PPL \\
\midrule
Sinusoidal PE (baseline) & 2.104 & 8.20 \\
NoPE & 2.265 & 9.63 \\
\midrule
Gap & +0.161 & +1.43 \\
\bottomrule
\end{tabular}
\end{table}

\begin{figure}[h]
  \centering
  \includegraphics[width=0.85\textwidth]{plots/nope_ablation.png}
  \caption{Sinusoidal PE vs.\ NoPE validation loss.  NoPE converges more
    slowly and finishes $\approx0.16$ nats above the sinusoidal baseline,
    confirming that explicit positional information helps but is not
    catastrophically necessary.}
  \label{fig:pe}
\end{figure}

\subsection{Observations and Explanation}

The performance gap is moderate ($\approx 0.16$ nats at 2{,}500 steps) rather
than catastrophic.  This is somewhat counterintuitive but has a principled
explanation:

\begin{itemize}
  \item \textbf{The causal mask provides implicit order information.}  In a
        decoder-only Transformer, the causal (lower-triangular) attention mask
        ensures that token $i$ can attend only to tokens $j \leq i$.  This
        means the \emph{relative order} of tokens is preserved even without
        explicit position IDs: the model can infer that an early-in-context
        token came before a later-in-context token by which ones were masked
        out.  However, it cannot distinguish between, say, position 5 and
        position 50 in absolute terms.

  \item \textbf{TinyStories has short-range dependencies.}  Most
        grammatical relationships in children's stories operate within a span
        of a few tokens (subject-verb agreement, pronoun resolution within a
        sentence).  The sinusoidal embeddings help most when long-range
        absolute position matters; on this dataset the benefit is real but
        modest.

  \item \textbf{NoPE is not position-blind.}  The relative ordering
        information from the causal mask means the model still learns
        \emph{recency} (recent tokens are more informative than distant ones),
        even without absolute position signals.  Learned attention patterns
        effectively implement a soft recency bias.
\end{itemize}

\textbf{Conclusion:}  Sinusoidal positional embeddings provide a measurable
but not decisive advantage on this dataset and context length (256 tokens).
For longer contexts or tasks requiring precise absolute positioning (e.g.,
copying text, structured generation), the gap would likely be larger.
For decoder-only models with moderate context lengths and local dependencies,
NoPE (or relative PE schemes) may be acceptable alternatives.

% ─────────────────────────────────────────────────────────────────────────────
\section{Summary of Results}
% ─────────────────────────────────────────────────────────────────────────────

\begin{table}[h]
\centering
\caption{All experimental results}
\label{tab:summary}
\begin{tabular}{llrr}
\toprule
Experiment & Configuration & Best Val Loss & Best PPL \\
\midrule
Baseline Phase 1 & lr=$3\times10^{-4}$, 5{,}000 steps & 2.104 & 8.20 \\
Baseline + Warm Restart & lr=$5\times10^{-4}$, 2{,}000 extra steps & \textbf{1.845} & \textbf{6.33} \\
\midrule
No LayerNorm (high LR) & lr=$3\times10^{-4}$, 400 steps & 3.534 & 34.26 \\
No LayerNorm (low LR) & lr=$3\times10^{-5}$, 1{,}500 steps & 3.948 & 51.85 \\
\midrule
NoPE & lr=$3\times10^{-4}$, 2{,}500 steps & 2.265 & 9.63 \\
\bottomrule
\end{tabular}
\end{table}

The best model (baseline with warm restart) achieves validation perplexity
$6.33$, comfortably below the target of $\leq 2.0$ validation loss.  The
ablations confirm that LayerNorm is essential for fast, stable convergence
(removing it raises perplexity by $4\times$--$8\times$), while sinusoidal
positional embeddings provide a useful but not critical benefit on short-context
language modeling tasks.

\end{document}
